\documentclass[UTF8,a4paper,11pt,fontset=fandol]{ctexart}
\usepackage[margin=2.2cm]{geometry}
\usepackage{amsmath,amssymb,booktabs}
\setlength{\parskip}{0.3em}
\linespread{1.1}

\begin{document}
\pagestyle{plain}

\begin{center}
  {\Large\bfseries 第三次上机作业：例 4.14 的 Markov 链极限分布}
\end{center}

\section{原题抄录}

编程重现教材第 130 页例 4.14，考察例 4.10 中的 Markov 链，
其转移概率矩阵为
\[
P=
\begin{pmatrix}
1/3 & 1/2 & 3/4\\
2/3 & 0   & 0\\
0   & 1/2 & 1/4
\end{pmatrix}.
\]
验证该 Markov 链不可约、非周期，计算其平稳分布，
并通过计算 $P^t$ 重现教材中的极限结果
\[
\lim_{t\to\infty}P^t
\approx
\begin{pmatrix}
0.4737 & 0.4737 & 0.4737\\
0.3158 & 0.3158 & 0.3158\\
0.2105 & 0.2105 & 0.2105
\end{pmatrix}.
\]

\section{代码的结构、算法说明}

\subsection*{（1）转移概率矩阵与分布的记号}

沿用教材的记号，状态空间为 $S=\{1,2,3\}$，
转移概率定义为
\[
p_{ij}=P[X_t=i\mid X_{t-1}=j].
\]
因此，矩阵第 $j$ 列表示从状态 $j$ 出发，
一步转移到各个状态的概率。
矩阵元素均非负，且每一列的元素之和为 1，即
\[
e^{\mathrm T}P=e^{\mathrm T},
\qquad e=(1,1,1)^{\mathrm T}.
\]

记 $t$ 时刻的概率分布为列向量 $\pi_t$，
则
\[
\pi_t=P\pi_{t-1}=P^t\pi_0.
\]
其中 $(P^t)_{ij}$ 表示经过 $t$ 步，
由状态 $j$ 转入状态 $i$ 的概率。

程序首先使用 \texttt{numpy.array} 定义矩阵 $P$，
并通过 \texttt{print\_array(name, A)}
统一输出向量和矩阵，各元素保留小数点后四位。
这一格式仅影响显示，不改变参与后续运算的数据精度。

\subsection*{（2）不可约与非周期的验证}

程序通过 \texttt{P2 = P @ P} 计算 $P^2$，
再用 \texttt{np.all(P2 > 0)}
判断其中的元素是否全部严格大于 0。

直接计算可得
\[
P^2=
\begin{pmatrix}
4/9 & 13/24 & 7/16\\
2/9 & 1/3   & 1/2\\
1/3 & 1/8   & 1/16
\end{pmatrix}>0.
\]
这里 $P^2>0$ 表示每个元素均严格大于 0。
因此，对任意状态 $i,j$，都有 $(P^2)_{ij}>0$，
即任意两个状态之间均存在两步转移路径，
从而该 Markov 链不可约。

非周期性还需要结合返回步数说明。
因为 $P$ 每列元素非负且总和为 1，
所以每列至少有一个正元素。于是
\[
(P^3)_{ij}
=\sum_{k=1}^{3}(P^2)_{ik}p_{kj}>0,
\]
从而 $P^3>0$。

按教材定义，状态 $i$ 的周期为
\[
d(i)=\gcd\{t\geq1:(P^t)_{ii}>0\}.
\]
由于 $(P^2)_{ii}>0$ 且 $(P^3)_{ii}>0$，
返回步数集合同时包含 2 和 3，其最大公约数为 1。
因此，对所有状态均有 $d(i)=1$，
该 Markov 链非周期。

由教材定理 4.20，该有限状态、齐次、不可约、
非周期的 Markov 链存在唯一的平稳分布 $\pi$，
且对任意初始概率分布 $\pi_0$，有
\[
\lim_{t\to\infty}P^t\pi_0=\pi.
\]

\subsection*{（3）平稳分布的计算}

平稳分布满足
\[
P\pi=\pi,\qquad
e^{\mathrm T}\pi=1,\qquad \pi_i\geq0.
\]
其中 $(P-I)\pi=0$ 的方程线性相关，
不能直接作为非奇异线性方程组求解。

程序先构造 \texttt{A = P - np.eye(n)}，
再将最后一行替换为全 1 向量，
同时将右端向量设为 $(0,0,1)^{\mathrm T}$。
因此，实际求解的方程组为
\[
\begin{pmatrix}
-2/3 & 1/2 & 3/4\\
2/3  & -1  & 0\\
1    & 1   & 1
\end{pmatrix}
\begin{pmatrix}
\pi_1\\\pi_2\\\pi_3
\end{pmatrix}
=
\begin{pmatrix}
0\\0\\1
\end{pmatrix}.
\]
调用 \texttt{numpy.linalg.solve(A, b)}
得到平稳分布的数值解。

由第二行方程可得
\[
\pi_2=\frac23\pi_1,
\]
代入第一行得
\[
\pi_3=\frac49\pi_1.
\]
结合归一化条件，得到精确结果
\[
\pi=\frac1{19}
\begin{pmatrix}
9\\6\\4
\end{pmatrix}.
\]
其各分量均为正，且总和为 1。

\subsection*{（4）极限矩阵与收敛误差}

分别取三个单位向量作为初始分布，
由定理 4.20 可知 $P^t$ 的每一列均收敛到 $\pi$。
因此
\[
\lim_{t\to\infty}P^t
=\pi e^{\mathrm T}
=\frac1{19}
\begin{pmatrix}
9&9&9\\
6&6&6\\
4&4&4
\end{pmatrix}.
\]
程序使用 \texttt{np.outer(pi, e)}
构造这一理论极限矩阵，保存为 \texttt{P\_limit}。

程序从 $P^0=I$ 出发，每一步执行
\[
P^{t}=P P^{t-1},
\]
即代码中的 \texttt{P\_t = P @ P\_t}，
依次计算至 $t=50$。

在 $t=1,2,5,10,20,50$ 时，
计算并输出误差
\[
E_t=
\max_{1\leq i,j\leq3}
\left|(P^t)_{ij}-\pi_i\right|.
\]
该指标是矩阵各元素误差绝对值的最大值，
不是按行求和定义的矩阵无穷范数。
程序使用
\texttt{np.max(np.abs(P\_t - P\_limit))}
计算这一指标，并以科学计数法保留六位小数输出。

程序固定进行 50 次矩阵乘法，
没有设置基于误差的提前停止条件。

\section{输出结果呈现}

运行程序，转移概率矩阵及其平方的输出为
\[
P=
\begin{pmatrix}
0.3333&0.5000&0.7500\\
0.6667&0.0000&0.0000\\
0.0000&0.5000&0.2500
\end{pmatrix},
\]
\[
P^2=
\begin{pmatrix}
0.4444&0.5417&0.4375\\
0.2222&0.3333&0.5000\\
0.3333&0.1250&0.0625
\end{pmatrix}.
\]
全部元素严格大于 0 的判断结果为
\begin{center}
\texttt{All entries of P\^{}2 are positive: True}
\end{center}

计算得到的平稳分布保留四位小数后为
\[
\pi\approx
\begin{pmatrix}
0.4737\\
0.3158\\
0.2105
\end{pmatrix}.
\]
这里程序使用一维 NumPy 数组保存 $\pi$，
因此控制台将其横向显示；
其数学含义仍为概率分布列向量。

理论极限矩阵与计算得到的 $P^{50}$
在程序的四位小数显示精度下相同：
\[
\pi e^{\mathrm T}\approx P^{50}\approx
\begin{pmatrix}
0.4737&0.4737&0.4737\\
0.3158&0.3158&0.3158\\
0.2105&0.2105&0.2105
\end{pmatrix}.
\]

各个选定步数对应的误差如下。
表中数据来自逐次左乘计算的一次运行，
最后几位可能受到浮点运算环境的影响。

\begin{center}
\begin{tabular}{cr}
\toprule
步数 $t$ & 误差 $E_t$\\
\midrule
1  & \texttt{3.508772e-01}\\
2  & \texttt{1.842105e-01}\\
5  & \texttt{1.262488e-02}\\
10 & \texttt{1.079195e-04}\\
20 & \texttt{1.871771e-08}\\
50 & \texttt{6.661338e-16}\\
\bottomrule
\end{tabular}
\end{center}

\section{对输出结果的分析和结论}

\subsection*{（1）定理条件的验证}

计算结果表明 $P^2$ 的每个元素均严格大于 0，
因此任意状态之间都可以经过两步相互到达，
该 Markov 链不可约。

同时，由 $P^2>0$ 可推出 $P^3>0$，
每个状态均可以在第 2 步和第 3 步以正概率返回自身，
所以周期为 1。
这里非周期性的依据是同时存在互素的返回步数；
仅知道 $(P^2)_{ii}>0$，
并不足以排除周期为 2 的情况。

因此，该 Markov 链满足教材定理 4.20 的条件，
平稳分布唯一，且任意初始概率分布均收敛到该平稳分布。

\subsection*{（2）平稳分布与教材结果的比较}

程序计算的平稳分布与精确结果
\[
\pi=\frac1{19}(9,6,4)^{\mathrm T}
\]
一致到浮点计算精度。
保留四位小数后为
$(0.4737,0.3158,0.2105)^{\mathrm T}$，
与教材例 4.14 中极限矩阵的每一列一致。

平稳分布的含义是：
无论初始概率分布如何，
经过足够多步转移后，
处于状态 1、2、3 的概率分别趋于
$9/19$、$6/19$ 和 $4/19$。
其中状态 1 的极限概率最大，状态 3 的极限概率最小。

\subsection*{（3）数值收敛过程}

在本次选取的步数上，误差随步数增加明显减小。
$t=5$ 时误差约为 $1.26\times10^{-2}$；
$t=10$ 时约为 $1.08\times10^{-4}$；
$t=20$ 时已减小到约 $1.87\times10^{-8}$。

到 $t=50$ 时，误差约为 $10^{-16}$，
已经接近双精度浮点数的舍入误差尺度。
这里的误差同时包含逐次矩阵乘法、
平稳分布求解以及极限矩阵构造中的浮点误差，
不能解释为精确算术下的全部截断误差。

此外，程序输出的矩阵只保留四位小数。
显示结果相同不意味着两个浮点矩阵逐元素完全相等，
因此需要结合误差表判断收敛程度。

\subsection*{（4）结论}

本实验通过计算 $P^2$ 并结合返回步数的性质，
验证了该 Markov 链不可约、非周期；
通过求解带归一化条件的线性方程组，
得到了唯一平稳分布
\[
\pi=\frac1{19}(9,6,4)^{\mathrm T}.
\]

逐次计算 $P^t$ 的结果表明，
矩阵各列逐渐趋于同一个平稳分布，
$P^{50}$ 在四位小数显示精度下与教材给出的极限矩阵一致，
完成了例 4.14 的数值重现。

\end{document}